\chapter{Phenomena}\label{ch:phenomena}

The first three chapters built a language and stated the rules of quantum mechanics in it. This chapter uses them. Each section takes a physical phenomenon, most of them famous, sets up the simplest model that captures it, and explains what happens with the tools developed so far. Three of them are the puzzles of the Introduction: interference in the double slit (Section~\ref{sec:interference}), the stability of atoms and their spectral lines (Section~\ref{sec:levels}), and black-body radiation (Section~\ref{sec:blackbody}). The other two, the Stern--Gerlach experiment with magnetic resonance (Section~\ref{sec:spin}) and the impossibility of cloning with quantum teleportation (Section~\ref{sec:teleportation}), show the postulates at work in the laboratory and in quantum information.

The models are deliberately simple: two-level systems, a particle in a box, a single mode of light. Simple models are not a weakness. They isolate the mechanism, and they are exactly what is used in practice to design qubits.

\begin{objectives}[In this chapter]
  You will explain interference and its destruction by which-path information; compute energy levels of a particle in a box and estimate the size and energy of the hydrogen atom; describe the Stern--Gerlach experiment, spin precession and magnetic resonance; derive Planck's law from the thermal state of a mode of light; and prove that unknown quantum states cannot be copied, but can be teleported.
\end{objectives}

\begin{whereitsits}
  This chapter uses all three previous ones: the postulates of Chapter~\ref{ch:postulates}, the operators and tensor products of Chapter~\ref{ch:tools}, and the density operators and entropies of Chapters~\ref{ch:postulates} and~\ref{ch:info}. Each section starts with a ``Recall'' note listing exactly what it uses.
\end{whereitsits}

%======================================================================
\section{Interference and which-path information}\label{sec:interference}
%======================================================================

\begin{objectives}
  You will model a two-path interferometer with a qubit, compute its interference fringes, see why ``bullets'' would show none, and show that recording which path the particle took destroys the fringes, by an amount fixed by an inner product.
\end{objectives}

\begin{recall}
  Relative phases and Example~\ref{ex:relative-phase} (Section~\ref{sec:complex}); the Hadamard gate (Section~\ref{sec:state-evolution}); measurement probabilities (Section~\ref{sec:measurement}); tensor products (Section~\ref{sec:tensor}); density operators and the partial trace (Section~\ref{sec:density}).
\end{recall}

The double slit of the Introduction (Figure~\ref{fig:slit}) has a continuum of points on its screen. Its essential mechanism, however, involves just two alternatives, the two slits, and a phase difference between them that changes along the screen. We therefore study the cleanest version of it: an interferometer with two paths and a phase that we control.

\subsection{The Mach--Zehnder interferometer}

A particle (a photon, say) enters a beam splitter, which sends it into a superposition of two paths, the upper one and the lower one. A phase shifter on the upper path adds a phase $\varphi$. The paths are recombined at a second beam splitter, and two detectors $D_0$ and $D_1$ record where the particle leaves (Figure~\ref{fig:mz}).

\begin{figure}[htb]
  \centering
  \begin{tikzpicture}[>=stealth,scale=1.15]
    % paths
    \draw[navy,thick] (-1.6,0) -- (0,0);
    \draw[navy,thick] (0,0) -- (3.5,0) -- (3.5,2.2);
    \draw[navy,thick] (0,0) -- (0,2.2) -- (3.5,2.2);
    \draw[navy,thick,->] (3.5,2.2) -- (5,2.2);
    \draw[navy,thick,->] (3.5,2.2) -- (3.5,3.6);
    % beam splitters
    \draw[ink,very thick] (-0.3,-0.3) -- (0.3,0.3);
    \draw[ink,very thick] (3.2,1.9) -- (3.8,2.5);
    % mirrors
    \draw[ink,line width=2.5pt] (-0.3,1.9) -- (0.3,2.5);
    \draw[ink,line width=2.5pt] (3.2,-0.3) -- (3.8,0.3);
    % phase shifter
    \draw[fill=figacc!25,draw=figacc] (1.45,1.95) rectangle (2.05,2.45);
    \node[figacc,font=\small] at (1.75,2.75) {$\varphi$};
    % labels
    \node[font=\footnotesize,anchor=east] at (-1.6,0) {source};
    \node[font=\footnotesize,anchor=north] at (0,-0.35) {BS$_1$};
    \node[font=\footnotesize,anchor=west] at (3.85,1.9) {BS$_2$};
    \fill[ink] (5.1,2.0) rectangle (5.4,2.4); \node[font=\footnotesize,anchor=west] at (5.45,2.2) {$D_0$};
    \fill[ink] (3.3,3.65) rectangle (3.7,3.95); \node[font=\footnotesize,anchor=south] at (3.5,3.95) {$D_1$};
    \node[font=\footnotesize,navy] at (1.75,-0.3) {lower path $\ket1$};
    \node[font=\footnotesize,navy,anchor=east] at (-0.1,1.1) {upper path $\ket0$};
  \end{tikzpicture}
  \caption{A Mach--Zehnder interferometer. The first beam splitter (BS$_1$) creates a superposition of the two paths, a phase shifter adds the phase $\varphi$ to one of them, and the second beam splitter (BS$_2$) recombines them in front of the detectors $D_0$, $D_1$.}
  \label{fig:mz}
\end{figure}

\paragraph{The model.}
The only thing that matters about the particle here is which path it is on, so its state is a qubit:
\begin{steps}
  \item $\ket0$ means ``upper path'' and $\ket1$ means ``lower path'';
  \item each beam splitter acts as the Hadamard gate $H$, sending $\ket0\mapsto\ket+$ and $\ket1\mapsto\ket-$ (a real beam splitter differs from $H$ by phases that do not change any result below);
  \item the phase shifter acts as $P_\varphi=\diag(1,e^{\ii\varphi})$ (it multiplies the amplitude of the upper path by $1$ and of the lower one by $e^{\ii\varphi}$; only the \emph{relative} phase matters);
  \item detector $D_0$ clicks when the final state is found in $\ket0$, and $D_1$ when it is found in $\ket1$: a measurement in the computational basis.
\end{steps}

\begin{proposition}[Interference fringes]\label{prop:mz}
  A particle entering in the state $\ket0$ reaches $D_0$ with probability $\cos^2(\varphi/2)$ and $D_1$ with probability $\sin^2(\varphi/2)$.
\end{proposition}
\begin{proof}
  \begin{steps}
    \item \emph{First beam splitter.} $H\ket0=\tfrac1{\sqrt2}(\ket0+\ket1)$: equal amplitudes for both paths.
    \item \emph{Phase shifter.} $P_\varphi$ gives $\tfrac1{\sqrt2}(\ket0+e^{\ii\varphi}\ket1)$.
    \item \emph{Second beam splitter.} Using $H\ket0=\ket+$ and $H\ket1=\ket-$, we get $\tfrac1{\sqrt2}(\ket++e^{\ii\varphi}\ket-)=\tfrac12\bigl((1+e^{\ii\varphi})\ket0+(1-e^{\ii\varphi})\ket1\bigr)$.
    \item \emph{Born rule.} $p(D_0)=\tfrac14\abs{1+e^{\ii\varphi}}^2=\tfrac14(2+2\cos\varphi)=\cos^2\tfrac\varphi2$, and similarly $p(D_1)=\tfrac14\abs{1-e^{\ii\varphi}}^2=\sin^2\tfrac\varphi2$. \qedhere
  \end{steps}
\end{proof}

As $\varphi$ varies, the probability of $D_0$ swings between $1$ and $0$. This is interference: each detector receives the \emph{sum} of the amplitudes of the two paths, $\tfrac12$ and $\tfrac12e^{\ii\varphi}$ for $D_0$, and these add constructively or cancel depending on their relative phase. It is the computation of Example~\ref{ex:relative-phase} in a laboratory. In the double slit, the role of $\varphi$ is played by the path difference, which changes from point to point of the screen; the bright and dark bands are the points where $\varphi$ is a multiple of $2\pi$ or an odd multiple of $\pi$.

\begin{example}[One particle at a time]\label{ex:mz-counts}
  Ten thousand particles are sent, one at a time, with $\varphi=\pi/3$. How many clicks of each detector do we expect, and what does a single particle do?
\end{example}
\begin{solution}
  \begin{steps}
    \item \emph{Probabilities.} $\cos^2(\pi/6)=\tfrac34$ and $\sin^2(\pi/6)=\tfrac14$.
    \item \emph{Counts.} About $7500$ clicks in $D_0$ and $2500$ in $D_1$; by the law of large numbers (Theorem~\ref{thm:wlln}) the fractions approach $\tfrac34,\tfrac14$ as the number of particles grows.
    \item \emph{A single particle} triggers exactly one detector, at random. The fringes are a statement about probabilities: they appear in the statistics of many particles, each of which interferes only with itself.
  \end{steps}
\end{solution}

\subsection{Why bullets show no fringes}

Suppose instead that each particle took a definite path, upper or lower with probability $\tfrac12$ each, as a bullet would. Then after the first beam splitter the path would not be in the superposition $\ket+$ but in the \emph{mixture} of $\ket0$ and $\ket1$, with density operator $\rho=\tfrac12\ket0\!\bra0+\tfrac12\ket1\!\bra1=\tfrac12\Id$.

\begin{example}[The classical mixture]\label{ex:mz-mixture}
  Compute the detection probabilities for the mixture $\rho=\tfrac12\Id$ entering the phase shifter.
\end{example}
\begin{solution}
  \begin{steps}
    \item \emph{Evolve} (Proposition~\ref{prop:postulates-rho}): $\rho\mapsto HP_\varphi\,\rho\,P_\varphi^\dagger H^\dagger=\tfrac12HP_\varphi P_\varphi^\dagger H^\dagger=\tfrac12\Id$, because the unitaries cancel against each other.
    \item \emph{Probabilities.} $p(D_0)=\bra0\tfrac12\Id\ket0=\tfrac12$ and $p(D_1)=\tfrac12$, for every $\varphi$.
  \end{steps}
  No fringes: probabilities add, not amplitudes. The superposition $\ket+$ and the mixture $\tfrac12\Id$ give the same statistics for each path separately, but only the superposition carries the relative phase between them (Example~\ref{ex:mixed-vs-superposition}).
\end{solution}

\subsection{Which-path information destroys interference}

The most striking feature of the double slit is that the fringes disappear if we find out which slit each particle went through. We can model this precisely. Put a second system, a \emph{marker} $M$, next to the upper and lower paths, so that it records the path: if the particle passes along path $j$, the marker goes to a state $\ket{m_j}$. Starting from the particle in $\ket+$ (after the first beam splitter) and the phase shifter, the joint state is
\[
  \ket\Psi=\tfrac1{\sqrt2}\bigl(\ket0\ket{m_0}+e^{\ii\varphi}\ket1\ket{m_1}\bigr).
\]
The marker could be anything: a photon scattered off the particle, an atom in a slit, a detector. The question is what happens at $D_0$ and $D_1$ when we do not look at the marker.

\begin{theorem}[Fringe visibility]\label{thm:visibility}
  Let $\gamma=\braket{m_0}{m_1}$. After the second beam splitter, $p(D_0)=\tfrac12\bigl(1+\abs\gamma\cos(\varphi+\arg\gamma)\bigr)$. The fringes therefore have visibility $\abs\gamma$: they are perfect when the marker states coincide, and absent when they are orthogonal.
\end{theorem}
\begin{proof}
  \begin{steps}
    \item \emph{Reduce to the particle.} Only the particle reaches the detectors, so its statistics are given by its reduced density operator $\rho=\Tr_M\ket\Psi\!\bra\Psi$ (Proposition~\ref{prop:reduced}).
    \item \emph{Compute it.} $\ket\Psi\!\bra\Psi$ has four terms. By the definition of the partial trace, $\Tr_M\bigl(\ket a\!\bra b\otimes\ket{m}\!\bra{m'}\bigr)=\braket{m'}{m}\ket a\!\bra b$, so
      \[
        \rho=\tfrac12\Bigl(\ket0\!\bra0+\ket1\!\bra1+e^{-\ii\varphi}\,\overline{\gamma}\,\ket0\!\bra1+e^{\ii\varphi}\gamma\,\ket1\!\bra0\Bigr).
      \]
    \item \emph{Second beam splitter and detection.} $p(D_0)=\bra0H\rho H\ket0=\bra+\rho\ket+=\tfrac12(\rho_{00}+\rho_{11}+\rho_{01}+\rho_{10})$.
    \item \emph{Substitute.} $\rho_{00}+\rho_{11}=1$ and $\rho_{01}+\rho_{10}=\tfrac12\bigl(e^{-\ii\varphi}\bar\gamma+e^{\ii\varphi}\gamma\bigr)=\Re(e^{\ii\varphi}\gamma)=\abs\gamma\cos(\varphi+\arg\gamma)$. \qedhere
  \end{steps}
\end{proof}

Two extreme cases make the meaning clear.
\begin{steps}
  \item If the marker does not record anything, $\ket{m_0}=\ket{m_1}$, then $\gamma=1$ and we recover Proposition~\ref{prop:mz}: full fringes.
  \item If the marker records the path perfectly, $\ket{m_0}\perp\ket{m_1}$, then $\gamma=0$, $\rho=\tfrac12\Id$, and $p(D_0)=\tfrac12$ for every $\varphi$: no fringes at all, exactly as for bullets. The entanglement with the marker has turned the superposition into a mixture.
\end{steps}
In between, the marker states are only partly distinguishable (Theorem~\ref{thm:no-distinguish}), and the fringes are only partly visible. Nothing in the argument requires anybody to \emph{read} the marker: it is enough that the information exists somewhere.

\begin{example}[Partial which-path information]\label{ex:partial-whichpath}
  Let $\ket{m_0}=\ket0$ and $\ket{m_1}=\cos\theta\ket0+\sin\theta\ket1$. Find the visibility of the fringes and the maximum of $p(D_0)$ as $\varphi$ varies, for $\theta=\pi/4$.
\end{example}
\begin{solution}
  \begin{steps}
    \item \emph{Overlap.} $\gamma=\braket{m_0}{m_1}=\cos\theta$, real and non-negative, so $\arg\gamma=0$.
    \item \emph{Visibility.} $\abs\gamma=\cos\tfrac\pi4=\tfrac1{\sqrt2}\approx0.71$.
    \item \emph{Range of $p(D_0)$.} By Theorem~\ref{thm:visibility}, $p(D_0)=\tfrac12\bigl(1+0.71\cos\varphi\bigr)$, which oscillates between $\tfrac12(1-0.71)\approx0.15$ and $\tfrac12(1+0.71)\approx0.85$, instead of between $0$ and $1$.
  \end{steps}
\end{solution}

\begin{lookahead}
  This is \emph{decoherence}. In a real experiment the role of the marker is played by the environment: stray photons, air molecules, the thermal vibrations of the apparatus. Each interaction leaves a little record of the path, makes $\abs\gamma$ smaller, and fades the fringes. It is why large objects never show interference in everyday life, and it is the main enemy of quantum computers, whose qubits must keep their relative phases for as long as a computation lasts.
\end{lookahead}

\begin{exerciseset}
  \exercise{1} Compute the detection probabilities of the interferometer for $\varphi=0,\ \pi/2,\ \pi$, and check them against Proposition~\ref{prop:mz}.

  \exercise{2} Replace both beam splitters by the symmetric beam splitter $B=\tfrac1{\sqrt2}\begin{pmatrix}1&\ii\\\ii&1\end{pmatrix}$. Show that $B$ is unitary and that the detection probabilities are still a constant plus a multiple of $\cos\varphi$. Which detector receives the particle when $\varphi=0$?

  \exercise{2} Block the lower path just after the first beam splitter (the particle is absorbed if it takes it). With what probability does a particle reach $D_0$, and does this depend on $\varphi$? Describe the state of the particle that is not absorbed.

  \exercise{2} Show that if the marker states are $\ket{m_0}=\ket0$ and $\ket{m_1}=e^{\ii\alpha}\ket0$, the fringes are fully visible but shifted. By how much?

  \exercise{3} The \emph{distinguishability} of the marker states is $D=\sqrt{1-\abs\gamma^2}$ (it is the largest difference between the probabilities that any measurement assigns to the two states). With the visibility $V=\abs\gamma$ of Theorem~\ref{thm:visibility}, show that $V^2+D^2=1$, and interpret: the more which-path information exists, the weaker the fringes.
\end{exerciseset}

%======================================================================
\section{Energy levels and the stability of atoms}\label{sec:levels}
%======================================================================

\begin{objectives}
  You will see why the energy of a confined particle comes in discrete levels, why a system in a single level does not change while a superposition of levels oscillates at the Bohr frequencies, how this explains sharp spectral lines, and why the uncertainty principle stops the electron of a hydrogen atom from falling into the nucleus.
\end{objectives}

\begin{recall}
  The spectral theorem and the functional calculus (Section~\ref{sec:operators}); the Schr\"odinger equation and energy eigenstates (Proposition~\ref{prop:schrodinger}); the uncertainty principle (Theorem~\ref{thm:uncertainty}) and the Gaussian wave packet (Example~\ref{ex:gaussian-uncertainty}); wave functions and the Born rule (Example~\ref{ex:box}).
\end{recall}

The second and third puzzles of the Introduction are about energy. Hydrogen emits light only at a few sharp frequencies (Figure~\ref{fig:levels}), and the atom does not collapse although classical physics says it must. Both are explained by one fact of the theory: the possible energies of a system are the eigenvalues of its Hamiltonian.

\subsection{Stationary states and Bohr frequencies}

Let $H=\sum_nE_n\ketbra{E_n}{E_n}$ be the spectral decomposition of the Hamiltonian of a system (Theorem~\ref{thm:spectral}), with $E_1<E_2<\cdots$ and, for simplicity, one eigenvector per eigenvalue. We saw after Proposition~\ref{prop:schrodinger} that a system in one energy eigenstate only picks up a global phase, $e^{-\ii Ht/\hbar}\ket{E_n}=e^{-\ii E_nt/\hbar}\ket{E_n}$, so nothing observable about it ever changes: it is a \term{stationary state}. A superposition of levels behaves differently.

\begin{proposition}[Only Bohr frequencies appear]\label{prop:bohr}
  Let $\ket{\psi(0)}=\sum_nc_n\ket{E_n}$ and let $A$ be any observable, with matrix elements $A_{mn}=\bra{E_m}A\ket{E_n}$. Then
  \[
    \langle A\rangle_{\psi(t)}=\sum_{m,n}\overline{c_m}\,c_n\,A_{mn}\,e^{\ii(E_m-E_n)t/\hbar}.
  \]
  The average of every observable is a combination of oscillations at the angular frequencies $\omega_{mn}=(E_m-E_n)/\hbar$, the \term{Bohr frequencies}, and of no others.
\end{proposition}
\begin{proof}
  \begin{steps}
    \item \emph{Evolve} (functional calculus): $\ket{\psi(t)}=e^{-\ii Ht/\hbar}\ket{\psi(0)}=\sum_nc_ne^{-\ii E_nt/\hbar}\ket{E_n}$.
    \item \emph{The bra} is $\bra{\psi(t)}=\sum_m\overline{c_m}e^{\ii E_mt/\hbar}\bra{E_m}$ (conjugate-linearity).
    \item \emph{Sandwich $A$:} $\bra{\psi(t)}A\ket{\psi(t)}=\sum_{m,n}\overline{c_m}c_n\,e^{\ii(E_m-E_n)t/\hbar}\bra{E_m}A\ket{E_n}$. \qedhere
  \end{steps}
\end{proof}

The diagonal terms $m=n$ are constant; only the cross terms move, and each one moves at the frequency of a \emph{difference} of two levels. This is the relative phase at work again: between two levels it grows like $(E_m-E_n)t/\hbar$, just as the relative phase between the two paths of an interferometer grew with $\varphi$.

\paragraph{Spectral lines.}
An electron is a charge. If its state is a superposition of two levels, its average position oscillates at a Bohr frequency, and an oscillating charge radiates light of that frequency, just as an antenna does. Hence the atom emits, and absorbs, light only at the frequencies
\[
  \nu_{mn}=\frac{\omega_{mn}}{2\pi}=\frac{E_m-E_n}{h},\qquad\text{that is,}\qquad h\nu_{mn}=E_m-E_n .
\]
This is Bohr's rule: the energy $h\nu$ of the light emitted equals the energy lost by the atom. A discrete set of levels gives a discrete set of lines. (A full theory of emission treats the light as a quantum system too, and Section~\ref{sec:blackbody} takes a first step in that direction; the frequencies, however, are already fixed by Proposition~\ref{prop:bohr}.)

\begin{example}[The red line of hydrogen]\label{ex:hydrogen-line}
  The levels of hydrogen are $E_n=-13.6\,\mathrm{eV}/n^2$. Find the energy and the wavelength of the light emitted when the electron drops from $n=3$ to $n=2$, using $E=hc/\lambda$ with $hc\approx1240\ \mathrm{eV\,nm}$.
\end{example}
\begin{solution}
  \begin{steps}
    \item \emph{The levels.} $E_3=-13.6/9\approx-1.51$ eV and $E_2=-13.6/4=-3.40$ eV.
    \item \emph{The energy of the light.} $h\nu=E_3-E_2\approx1.89$ eV.
    \item \emph{The wavelength.} $\lambda=hc/(h\nu)\approx1240/1.89\approx656$ nm.
  \end{steps}
  This is the bright red line of hydrogen, and the drops from $n=4,5,\dots$ to $n=2$ give the other visible lines. One formula for the levels explains the whole spectrum; what remains is to understand why the levels are discrete at all.
\end{solution}

\subsection{The particle in a box}

The simplest system with discrete levels is a particle of mass $m$ that moves freely on a segment $[0,L]$ and cannot leave it (a ``box''). Its state is a wave function $\psi(x)$ with $\psi(0)=\psi(L)=0$, and its Hamiltonian is the kinetic energy
\[
  H=\frac{\hat p^2}{2m}=-\frac{\hbar^2}{2m}\frac{d^2}{dx^2},\qquad\hat p=-\ii\hbar\frac{d}{dx}
\]
(Exercise~\ref{ex:no-xp} introduced $\hat p$). The space of wave functions is infinite-dimensional, which Chapter~\ref{ch:tools} did not cover; but we shall see that $H$ has an orthonormal basis of eigenvectors, exactly as the spectral theorem would lead us to expect, and that is all we need.

\begin{theorem}[Levels of a particle in a box]\label{thm:box}
  The eigenvalue problem $H\varphi=E\varphi$ with $\varphi(0)=\varphi(L)=0$ and $\varphi\neq0$ has the solutions
  \[
    E_n=\frac{n^2\pi^2\hbar^2}{2mL^2},\qquad\varphi_n(x)=\sqrt{\frac2L}\,\sin\frac{n\pi x}{L},\qquad n=1,2,3,\dots,
  \]
  and the functions $\varphi_n$ are orthonormal: $\int_0^L\varphi_m\varphi_n\dd x=\delta_{mn}$.
\end{theorem}
\begin{proof}
  \begin{steps}
    \item \emph{Rewrite the equation.} $H\varphi=E\varphi$ means $\varphi''=-k^2\varphi$ with $k^2=2mE/\hbar^2$.
    \item \emph{The energy is positive.} Multiply $-\frac{\hbar^2}{2m}\varphi''=E\varphi$ by $\overline\varphi$ and integrate by parts, using $\varphi(0)=\varphi(L)=0$: $E\int_0^L\abs\varphi^2\dd x=\frac{\hbar^2}{2m}\int_0^L\abs{\varphi'}^2\dd x$. The right side is $>0$ (if $\varphi'=0$ then $\varphi$ is constant, hence $0$), so $E>0$ and $k>0$.
    \item \emph{General solution.} The solutions of $\varphi''=-k^2\varphi$ are $\varphi(x)=a\sin kx+b\cos kx$.
    \item \emph{Left wall.} $\varphi(0)=b=0$.
    \item \emph{Right wall.} $\varphi(L)=a\sin kL=0$ with $a\neq0$ forces $kL=n\pi$ for an integer $n\ge1$. Then $E=\hbar^2k^2/2m=n^2\pi^2\hbar^2/2mL^2$.
    \item \emph{Normalise.} $\int_0^L\sin^2\frac{n\pi x}{L}\dd x=\frac L2$, so $a=\sqrt{2/L}$.
    \item \emph{Orthogonality.} For $m\neq n$, $2\sin\alpha\sin\beta=\cos(\alpha-\beta)-\cos(\alpha+\beta)$ gives $\int_0^L\varphi_m\varphi_n\dd x=\frac1L\int_0^L\bigl(\cos\frac{(m-n)\pi x}L-\cos\frac{(m+n)\pi x}L\bigr)\dd x=0$. \qedhere
  \end{steps}
\end{proof}

\begin{figure}[htb]
  \centering
  \begin{tikzpicture}
    \begin{axis}[width=0.62\linewidth, height=6.2cm, xmin=-0.08, xmax=1.08, ymin=-0.3, ymax=10.6,
        axis lines=none, clip=false]
      \draw[ink, very thick] (axis cs:0,-0.2) -- (axis cs:0,10.4);
      \draw[ink, very thick] (axis cs:1,-0.2) -- (axis cs:1,10.4);
      \draw[ink!60, dashed] (axis cs:0,1) -- (axis cs:1,1);
      \draw[ink!60, dashed] (axis cs:0,4) -- (axis cs:1,4);
      \draw[ink!60, dashed] (axis cs:0,9) -- (axis cs:1,9);
      \node[anchor=west,font=\footnotesize] at (axis cs:1.03,1) {$E_1$};
      \node[anchor=west,font=\footnotesize] at (axis cs:1.03,4) {$E_2=4E_1$};
      \node[anchor=west,font=\footnotesize] at (axis cs:1.03,9) {$E_3=9E_1$};
      \addplot[navy, thick, samples=80, domain=0:1] {1+0.9*sin(deg(pi*x))};
      \addplot[navy, thick, samples=80, domain=0:1] {4+0.9*sin(deg(2*pi*x))};
      \addplot[navy, thick, samples=80, domain=0:1] {9+0.9*sin(deg(3*pi*x))};
      \node[font=\footnotesize,anchor=north] at (axis cs:0,-0.25) {$0$};
      \node[font=\footnotesize,anchor=north] at (axis cs:1,-0.25) {$L$};
    \end{axis}
  \end{tikzpicture}
  \caption{The three lowest levels of a particle in a box. Each eigenfunction $\varphi_n$ is drawn at the height of its energy: it fits $n$ half-wavelengths between the walls, and the energies grow like $n^2$.}
  \label{fig:box}
\end{figure}

The discreteness comes from the boundary conditions. A wave that must vanish at both walls has to fit a whole number of half-wavelengths into the box, as the vibrations of a guitar string do (Figure~\ref{fig:box}); each allowed wavelength $2L/n$ gives one allowed momentum $\hbar k=n\pi\hbar/L$ and one allowed energy. A free particle, with no walls, has no such restriction, and its energies form a continuum (Example~\ref{ex:gaussian-uncertainty} used such free wave functions).

Two further features deserve comment.
\begin{steps}
  \item \emph{The lowest energy is not zero.} $E_1=\pi^2\hbar^2/2mL^2>0$: a confined particle can never be at rest. This \term{zero-point energy} is forced by the uncertainty principle. Confinement to $[0,L]$ makes $\Delta x$ at most of order $L$, so $\Delta p\ge\hbar/2\Delta x$ is at least of order $\hbar/L$, and the average kinetic energy $\langle\hat p^2\rangle/2m\ge(\Delta p)^2/2m$ is at least of order $\hbar^2/mL^2$.
  \item \emph{Smaller boxes have more widely spaced levels.} The levels scale like $1/L^2$.
\end{steps}

\begin{example}[An electron in a nanometre box]\label{ex:nano-box}
  An electron ($m\approx9.11\times10^{-31}$ kg) is confined to a box of length $L=1$ nm. Find $E_1$, the energy of the transition $2\to1$, and the wavelength of the corresponding light. Use $\hbar\approx1.055\times10^{-34}$ J\,s and $1\ \mathrm{eV}\approx1.602\times10^{-19}$ J.
\end{example}
\begin{solution}
  \begin{steps}
    \item \emph{Ground level.} $E_1=\dfrac{\pi^2\hbar^2}{2mL^2}=\dfrac{(3.31\times10^{-34})^2}{2\cdot9.11\times10^{-31}\cdot10^{-18}}\approx6.0\times10^{-20}\ \mathrm J\approx0.38\ \mathrm{eV}$.
    \item \emph{Transition.} $E_2-E_1=(4-1)E_1=3E_1\approx1.13$ eV.
    \item \emph{Wavelength.} $\lambda\approx1240/1.13\approx1100$ nm, in the near infrared.
  \end{steps}
  Energies of an electronvolt for sizes of a nanometre: this is the scale of atoms and of chemistry. Nanocrystals called \emph{quantum dots} are boxes of this kind, and their colour is set by their size: smaller dots have more widely spaced levels and glow bluer.
\end{solution}

\begin{example}[A superposition in the box moves]\label{ex:box-moves}
  Let $\ket{\psi(0)}=\tfrac1{\sqrt2}(\varphi_1+\varphi_2)$. Show that the average position oscillates about the centre of the box at the Bohr frequency $\omega_{21}=(E_2-E_1)/\hbar$, and find the amplitude of the oscillation, given that $\int_0^Lx\,\varphi_1\varphi_2\dd x=-\tfrac{16L}{9\pi^2}$.
\end{example}
\begin{solution}
  \begin{steps}
    \item \emph{Diagonal terms.} $\int_0^Lx\,\varphi_n^2\dd x=\tfrac L2$ for every $n$, because $\varphi_n^2$ is symmetric about $x=\tfrac L2$.
    \item \emph{Apply Proposition~\ref{prop:bohr}} with $A=\hat x$, $c_1=c_2=\tfrac1{\sqrt2}$, and $x_{12}=x_{21}=-\tfrac{16L}{9\pi^2}$ (the eigenfunctions are real):
      \[
        \langle\hat x\rangle(t)=\tfrac12\cdot\tfrac L2+\tfrac12\cdot\tfrac L2+\tfrac12x_{12}\bigl(e^{-\ii\omega_{21}t}+e^{\ii\omega_{21}t}\bigr)=\tfrac L2-\tfrac{16L}{9\pi^2}\cos\omega_{21}t .
      \]
    \item \emph{Amplitude.} $\tfrac{16}{9\pi^2}\approx0.18$: the particle sloshes from side to side with an amplitude of about $18\%$ of the box.
  \end{steps}
  A charged particle doing this radiates at frequency $\omega_{21}/2\pi$, which is the mechanism behind the spectral lines described above. Each stationary state alone gives a constant $\langle\hat x\rangle=L/2$.
\end{solution}

\subsection{Why atoms do not collapse}

The hydrogen atom has the Hamiltonian
\[
  H=\frac{\hat p^2}{2m}-\frac{ke^2}{r},\qquad k=\frac1{4\pi\varepsilon_0},
\]
kinetic energy plus the Coulomb attraction to the nucleus at distance $r$. Classically the electron can lower its energy without limit by falling inwards, since $-ke^2/r\to-\infty$. Quantum mechanically, falling inwards means confining the electron to a smaller region, and confinement costs kinetic energy, as we have just seen. The two effects balance at a definite size.

\begin{example}[Size and energy of the hydrogen atom]\label{ex:hydrogen-estimate}
  Suppose that the electron is spread over a region of radius $r$ around the nucleus. Estimate its energy as a function of $r$, find the radius that minimises it, and evaluate the radius and the energy numerically.
\end{example}
\begin{solution}
  \begin{steps}
    \item \emph{Kinetic energy.} Confinement to size $r$ gives $\Delta x\sim r$, hence $\Delta p\sim\hbar/r$ by the uncertainty principle, and a kinetic energy of order $\hbar^2/2mr^2$ (we drop factors of order one).
    \item \emph{Potential energy.} Of order $-ke^2/r$.
    \item \emph{Total.} $E(r)\approx\dfrac{\hbar^2}{2mr^2}-\dfrac{ke^2}{r}$. The first term wins at small $r$ and the second at large $r$, so $E(r)$ has a minimum.
    \item \emph{Minimise.} $E'(r)=-\dfrac{\hbar^2}{mr^3}+\dfrac{ke^2}{r^2}=0$ gives
      \[
        r=a_0=\frac{\hbar^2}{mke^2},\qquad E(a_0)=\frac{ke^2}{2a_0}-\frac{ke^2}{a_0}=-\frac{ke^2}{2a_0}.
      \]
    \item \emph{Numbers.} With $ke^2\approx2.31\times10^{-28}$ J\,m, $a_0\approx5.3\times10^{-11}$ m, about half an \aa ngstr\"om, and $E(a_0)\approx-2.18\times10^{-18}$ J $\approx-13.6$ eV.
  \end{steps}
  The estimate produces the size of the atom and, by a lucky choice of the factors we dropped, exactly the ground level $E_1=-13.6$ eV of Figure~\ref{fig:levels}. The exact solution of the eigenvalue problem for $H$ gives all the levels $E_n=-13.6\,\mathrm{eV}/n^2$; it is a longer computation, done in every textbook of quantum mechanics~\cite{griffiths2018}.
\end{solution}

So the atom is stable because its Hamiltonian is bounded below: there is a lowest eigenvalue, the ground level, and an atom in its ground state is in a stationary state that has nowhere lower to go. The uncertainty principle, which in Chapter~\ref{ch:postulates} looked like a limitation on what we can know, is here what keeps matter from collapsing.

\begin{lookahead}
  A qubit is, physically, two levels of some system, chosen so that they are well separated from all the others: two levels of an ion, of an atom, or of a superconducting circuit that behaves like an artificial atom. The Bohr frequency between them is the frequency of the light or microwaves used to control them, as the next section shows.
\end{lookahead}

\begin{exerciseset}
  \exercise{1} Find the energy and the wavelength of the light emitted in the transition $n=4\to n=2$ of hydrogen. Is it visible (roughly $400$--$700$ nm)?

  \exercise{1} Verify directly that $\varphi_2(x)=\sqrt{2/L}\,\sin(2\pi x/L)$ satisfies $H\varphi_2=E_2\varphi_2$ and the boundary conditions, and that $\int_0^L\varphi_1\varphi_2\dd x=0$.

  \exercise{2} For the ground state $\varphi_1$ of the box, compute $\langle\hat x\rangle$, $\Delta x$, $\langle\hat p\rangle$ and $\Delta p$. Show that $\Delta x\,\Delta p=\hbar\sqrt{\pi^2/12-1/2}\approx0.57\hbar$, and compare with Theorem~\ref{thm:uncertainty}.

  \exercise{2} Compute $\int_0^Lx\,\varphi_1\varphi_2\dd x=-\tfrac{16L}{9\pi^2}$, the value used in Example~\ref{ex:box-moves}.

  \exercise{2} Estimate the length of a box in which an electron emits red light ($650$ nm) in the transition $2\to1$.

  \exercise{2} Repeat the estimate of Example~\ref{ex:hydrogen-estimate} for an electron bound to a nucleus of charge $Ze$. Show that the radius scales like $1/Z$ and the energy like $Z^2$, and estimate the ground energy of the ion He$^+$ ($Z=2$).

  \exercise{3} Let the spectrum of $H$ be non-degenerate and $\ket{\psi(0)}=\sum_nc_n\ket{E_n}$. Show that the time average $\lim_{T\to\infty}\frac1T\int_0^T\langle A\rangle_{\psi(t)}\dd t$ equals $\sum_n\abs{c_n}^2A_{nn}$, so that over long times only the populations of the levels matter. (Assume the sum is finite.)
\end{exerciseset}

%======================================================================
\section{Spin: Stern--Gerlach, precession and resonance}\label{sec:spin}
%======================================================================

\begin{objectives}
  You will describe the Stern--Gerlach experiment and its sequential versions with a qubit, show that a spin in a magnetic field precesses, derive the Rabi formula for magnetic resonance (the way qubits are actually flipped in the laboratory), and see how frequent measurements can freeze the evolution of a system.
\end{objectives}

\begin{recall}
  The Pauli matrices and their products (Proposition~\ref{prop:pauli}); the Bloch sphere and the overlap formula (Proposition~\ref{prop:bloch}); the Schr\"odinger equation and Examples~\ref{ex:H-eq-X} and~\ref{ex:precession} (Section~\ref{sec:state-evolution}); projective measurements (Section~\ref{sec:measurement}).
\end{recall}

\subsection{The Stern--Gerlach experiment}

In 1922 Stern and Gerlach~\cite{gerlach1922} sent a beam of silver atoms through a magnetic field that was stronger near one pole of the magnet than near the other. Each atom behaves as a tiny magnet, and in such a field a magnet is pushed up or down according to how it is oriented. Classically the atoms come out of the oven oriented at random, so the beam should be smeared into a continuous band. Instead it split into exactly \emph{two} spots.

The quantity that is measured is the \term{spin} of the atom (in silver it is the spin of one unpaired electron) along the direction of the field. The experiment shows that it has only two possible values, which are $\pm\hbar/2$. This is the textbook example of a qubit.

\paragraph{The model.}
\begin{steps}
  \item The spin state is a qubit, $\ket\psi\in\C^2$; $\ket0$ is ``spin up along $z$'' and $\ket1$ is ``spin down along $z$''.
  \item The spin along a unit vector $\vec n\in\R^3$ is the observable $S_{\vec n}=\tfrac\hbar2\,\vec n\cdot\vec\sigma$. Its eigenvalues are $\pm\tfrac\hbar2$ (Exercise~\ref{ex:spin-eig}), with unit eigenvectors $\ket{\pm\vec n}$, the states whose Bloch vectors are $\pm\vec n$.
  \item A Stern--Gerlach apparatus oriented along $\vec n$ (we write SG$_{\vec n}$) measures $S_{\vec n}$: it sends the atom to the upper spot with probability $\abs{\braket{+\vec n}\psi}^2$, and the atom leaves in the state $\ket{+\vec n}$; similarly for the lower spot.
\end{steps}
By Proposition~\ref{prop:bloch}, if the atom has Bloch vector $\vec r$, the probability of the upper spot is
\begin{equation}\label{eq:sg-prob}
  p_+=\abs{\braket{+\vec n}\psi}^2=\tfrac12\bigl(1+\vec n\cdot\vec r\bigr).
\end{equation}
Atoms from an oven are in the maximally mixed state $\rho=\tfrac12\Id$, and give $p_+=\tfrac12$ along every direction: two spots of equal intensity, whatever the orientation of the magnet.

\begin{figure}[htb]
  \centering
  \begin{tikzpicture}[>=stealth,
      sg/.style={draw=navy,thick,fill=navy!8,minimum width=1.3cm,minimum height=0.9cm,font=\small},
      lab/.style={font=\footnotesize}]
    \node[draw=ink,thick,fill=ink!10,minimum width=1cm,minimum height=0.7cm,font=\footnotesize] (oven) at (0,0) {oven};
    \node[sg] (A) at (2.4,0) {SG$_z$};
    \node[sg] (B) at (5.4,0.45) {SG$_x$};
    \node[sg] (C) at (8.4,0.9) {SG$_z$};
    \draw[->,thick] (oven) -- (A);
    \draw[->,thick] (A.east) ++(0,0.2) -- node[lab,above] {$+$} (B.west);
    \draw[thick] (A.east) ++(0,-0.2) -- ++(0.6,-0.35) node[lab,right] {$-$, blocked};
    \draw[->,thick] (B.east) ++(0,0.2) -- node[lab,above] {$+$} (C.west);
    \draw[thick] (B.east) ++(0,-0.2) -- ++(0.6,-0.35) node[lab,right] {$-$, blocked};
    \draw[->,thick] (C.east) ++(0,0.2) -- ++(0.8,0.3) node[lab,right] {$+$: $\tfrac12$};
    \draw[->,thick] (C.east) ++(0,-0.2) -- ++(0.8,-0.3) node[lab,right] {$-$: $\tfrac12$};
  \end{tikzpicture}
  \caption{Three Stern--Gerlach apparatuses in a row. The first one selects spin up along $z$, the second one selects spin up along $x$, and the third one measures the spin along $z$ again: both results reappear, with probability $\tfrac12$ each.}
  \label{fig:sg}
\end{figure}

\begin{example}[Measurements in a row]\label{ex:sg-sequence}
  Atoms pass through the three apparatuses of Figure~\ref{fig:sg}. Of the atoms that leave the oven, what fraction reaches each final spot? What happens if the middle apparatus is removed?
\end{example}
\begin{solution}
  \begin{steps}
    \item \emph{First SG$_z$.} From $\rho=\tfrac12\Id$, half of the atoms go up; they leave in $\ket0$, Bloch vector $(0,0,1)$.
    \item \emph{SG$_x$.} By \eqref{eq:sg-prob} with $\vec n=(1,0,0)$, $p_+=\tfrac12(1+0)=\tfrac12$. The selected atoms leave in $\ket+$, Bloch vector $(1,0,0)$.
    \item \emph{Last SG$_z$.} Now $p_\pm=\tfrac12(1\pm0)=\tfrac12$.
    \item \emph{Fractions.} $\tfrac12\cdot\tfrac12\cdot\tfrac12=\tfrac18$ of the atoms reaches each final spot.
    \item \emph{Without the middle apparatus,} the atoms reach the last SG$_z$ in $\ket0$, and all of them go up: the fractions are $\tfrac12$ and $0$.
  \end{steps}
  The atoms selected as ``up along $z$'' are found ``down along $z$'' after an $x$ measurement. The measurement of $S_x$ has not revealed a property that was there; it has changed the state, from $\ket0$ to $\ket+$ (Postulate~\ref{post:measurement}), and erased the information about $S_z$. This is the content of the non-commutativity $\comm{\sx}{\sz}\neq0$, and of the uncertainty relation $\Delta S_x\,\Delta S_z\ge\tfrac12\abs{\langle\comm{S_x}{S_z}\rangle}$.
\end{solution}

\subsection{Precession in a magnetic field}

A spin in a constant magnetic field $\vec B$ has energy proportional to its component along the field. With the field along $z$ the Hamiltonian is
\[
  H=\frac{\hbar\omega_0}2\,\sz,
\]
where $\omega_0$, the \term{Larmor frequency}, is proportional to $\abs{\vec B}$. The two levels are $\ket0$ and $\ket1$, with energies $\pm\hbar\omega_0/2$, and their Bohr frequency is $\omega_0$. To see what the spin does, we follow the averages $\langle\sigma_j\rangle$, which form its Bloch vector (Proposition~\ref{prop:bloch}).

\begin{lemma}[Evolution of averages]\label{lem:heisenberg}
  If $\ket{\psi(t)}$ obeys the Schr\"odinger equation with Hamiltonian $H$ and $A$ is a fixed operator, then
  \[
    \frac{d}{dt}\langle A\rangle_{\psi(t)}=\frac\ii\hbar\,\bigl\langle\comm HA\bigr\rangle_{\psi(t)} .
  \]
\end{lemma}
\begin{proof}
  \begin{steps}
    \item \emph{The equation and its adjoint.} $\frac{d}{dt}\ket\psi=-\frac\ii\hbar H\ket\psi$, and taking adjoints, $\frac{d}{dt}\bra\psi=\frac\ii\hbar\bra\psi H$ (because $H=H^\dagger$).
    \item \emph{Product rule.} $\frac{d}{dt}\bra\psi A\ket\psi=\frac\ii\hbar\bra\psi HA\ket\psi-\frac\ii\hbar\bra\psi AH\ket\psi=\frac\ii\hbar\bra\psi\comm HA\ket\psi$. \qedhere
  \end{steps}
\end{proof}

In particular an observable that commutes with $H$ has a constant average: the energy itself is conserved.

\begin{proposition}[Larmor precession]\label{prop:larmor}
  For $H=\tfrac{\hbar\omega_0}2\sz$, the Bloch vector $\vec r(t)$ of the state rotates about the $z$-axis at angular velocity $\omega_0$:
  \[
    \begin{aligned}
      r_x(t)&=r_x(0)\cos\omega_0t-r_y(0)\sin\omega_0t,\\
      r_y(t)&=r_x(0)\sin\omega_0t+r_y(0)\cos\omega_0t,\\
      r_z(t)&=r_z(0).
    \end{aligned}
  \]
\end{proposition}
\begin{proof}
  \begin{steps}
    \item \emph{The commutators} (Proposition~\ref{prop:pauli}): $\comm\sz\sx=\sz\sx-\sx\sz=2\ii\sy$, $\comm\sz\sy=-2\ii\sx$, and $\comm\sz\sz=0$.
    \item \emph{Apply Lemma~\ref{lem:heisenberg}} to $A=\sigma_j$, with $\tfrac\ii\hbar\cdot\tfrac{\hbar\omega_0}{2}=\tfrac{\ii\omega_0}2$:
      \[
        \dot r_x=\tfrac{\ii\omega_0}2\cdot2\ii\,r_y=-\omega_0r_y,\qquad\dot r_y=\tfrac{\ii\omega_0}2\cdot(-2\ii)\,r_x=\omega_0r_x,\qquad\dot r_z=0 .
      \]
    \item \emph{Solve.} The functions in the statement satisfy these equations and the initial conditions, and the solution of a linear system of differential equations with given initial values is unique. \qedhere
  \end{steps}
\end{proof}

This is Example~\ref{ex:precession} again, now for every initial state: the spin turns around the field like a spinning top, at the Bohr frequency of its two levels. Nothing in the argument depends on the field being along $z$: for $H=\tfrac{\hbar\omega_0}2\,\vec n\cdot\vec\sigma$ the Bloch vector rotates about $\vec n$ (Exercise~\ref{ex:precess-n}).

\subsection{Magnetic resonance and Rabi oscillations}

How does one turn $\ket0$ into $\ket1$, that is, apply the gate $X$? Precession about $z$ cannot do it: it never changes $r_z$. The answer, discovered by Rabi in the 1930s for atomic beams and used today in magnetic resonance imaging and in quantum computers, is to add a weak field that rotates in the $xy$-plane at a frequency $\omega$ close to the Larmor frequency $\omega_0$:
\[
  H(t)=\frac{\hbar\omega_0}2\,\sz+\frac{\hbar\Omega}2\bigl(\cos(\omega t)\,\sx+\sin(\omega t)\,\sy\bigr),
\]
where the \term{Rabi frequency} $\Omega>0$ measures the strength of the rotating field and $\Delta=\omega_0-\omega$ is the \term{detuning}.

\begin{theorem}[Rabi formula]\label{thm:rabi}
  A spin that starts in $\ket0$ and evolves with $H(t)$ is found in $\ket1$ at time $t$ with probability
  \[
    p_1(t)=\frac{\Omega^2}{\Omega^2+\Delta^2}\,\sin^2\Bigl(\frac{\sqrt{\Omega^2+\Delta^2}}{2}\,t\Bigr).
  \]
\end{theorem}
\begin{proof}
  The trick is to look at the spin from a frame that rotates with the field.
  \begin{steps}
    \item \emph{Write the field as a rotated $\sx$.} Let $R(t)=e^{-\ii\omega t\sz/2}$. A direct computation with the matrices gives $R(t)\,\sx\,R(t)^\dagger=\cos(\omega t)\sx+\sin(\omega t)\sy$ (Exercise~\ref{ex:rotate-x}), and $R(t)$ commutes with $\sz$. Hence $H(t)=R(t)H_1R(t)^\dagger$ with the \emph{constant} operator $H_1=\tfrac{\hbar\omega_0}2\sz+\tfrac{\hbar\Omega}2\sx$.
    \item \emph{Change variables.} Put $\ket{\phi(t)}=R(t)^\dagger\ket{\psi(t)}$. Using $\frac{d}{dt}R^\dagger=\tfrac{\ii\omega}2\sz R^\dagger$ and the Schr\"odinger equation for $\ket\psi$,
      \[
        \ii\hbar\frac{d}{dt}\ket\phi=-\tfrac{\hbar\omega}2\sz\ket\phi+R^\dagger H\ket\psi=\Bigl(-\tfrac{\hbar\omega}2\sz+H_1\Bigr)\ket\phi=\tfrac\hbar2\bigl(\Delta\,\sz+\Omega\,\sx\bigr)\ket\phi .
      \]
      In the rotating frame the Hamiltonian $H'=\tfrac\hbar2(\Delta\sz+\Omega\sx)$ does not depend on time.
    \item \emph{Write $H'$ as a field along a unit vector.} With $\Omega_R=\sqrt{\Omega^2+\Delta^2}$ and $\vec m=(\Omega,0,\Delta)/\Omega_R$, $H'=\tfrac{\hbar\Omega_R}2\,\vec m\cdot\vec\sigma$. Since $(\vec m\cdot\vec\sigma)^2=\Id$, the exponential is $e^{-\ii H't/\hbar}=\cos\tfrac{\Omega_Rt}2\,\Id-\ii\sin\tfrac{\Omega_Rt}2\,\vec m\cdot\vec\sigma$ (Exercise~\ref{ex:rotation-formula}).
    \item \emph{The amplitude of $\ket1$.} $\ket{\phi(0)}=\ket0$, and $\bra1\vec m\cdot\vec\sigma\ket0=m_x\bra1\sx\ket0+m_z\bra1\sz\ket0=m_x$. So $\braket1{\phi(t)}=-\ii\,\tfrac{\Omega}{\Omega_R}\sin\tfrac{\Omega_Rt}2$.
    \item \emph{Back to the laboratory.} $R(t)$ is diagonal, so it changes the amplitude of $\ket1$ only by a phase: $\abs{\braket1{\psi(t)}}=\abs{\braket1{\phi(t)}}$. Squaring gives the formula. \qedhere
  \end{steps}
\end{proof}

\begin{figure}[htb]
  \centering
  \begin{tikzpicture}
    \begin{axis}[width=0.8\linewidth, height=5cm, xmin=0, xmax=12.566, ymin=0, ymax=1.15,
        xlabel={time, in units of $1/\Omega$}, ylabel={$p_1(t)$},
        xtick={0,3.1416,6.2832,9.4248,12.566}, xticklabels={$0$,$\pi$,$2\pi$,$3\pi$,$4\pi$},
        ytick={0,0.5,1}, axis lines=left,
        legend style={font=\footnotesize, at={(0.5,1.0)}, anchor=south, draw=none, legend columns=2},
        label style={font=\footnotesize}, tick label style={font=\footnotesize}]
      \addplot[navy, very thick, samples=200, domain=0:12.566] {sin(deg(x/2))^2};
      \addplot[figacc, very thick, dashed, samples=200, domain=0:12.566] {0.5*sin(deg(sqrt(2)*x/2))^2};
      \legend{on resonance ($\Delta=0$), detuned ($\Delta=\Omega$)}
    \end{axis}
  \end{tikzpicture}
  \caption{Rabi oscillations. On resonance the spin is flipped completely at $t=\pi/\Omega$; away from resonance the oscillation is faster but incomplete.}
  \label{fig:rabi}
\end{figure}

Two features stand out (Figure~\ref{fig:rabi}).
\begin{steps}
  \item \emph{Resonance.} The largest probability of a flip, $\Omega^2/(\Omega^2+\Delta^2)$, equals $1$ only when $\omega=\omega_0$, and falls off quickly once $\abs\Delta$ exceeds $\Omega$. A weak field has a large effect only at the Bohr frequency of the two levels. Scanning $\omega$ and watching for flips is therefore a very precise way of measuring $\omega_0$, which depends on the local magnetic field and on the chemical surroundings of the spin: this is the principle of nuclear magnetic resonance and of MRI scanners.
  \item \emph{Pulses are gates.} On resonance, $H'=\tfrac{\hbar\Omega}2\sx$, so a pulse of duration $t$ applies, in the rotating frame, the rotation $e^{-\ii\Omega t\sx/2}$ about the $x$-axis by the angle $\Omega t$. A ``$\pi$-pulse'' ($\Omega t=\pi$) gives $-\ii\sx$, the gate $X$ up to a global phase; a ``$\pi/2$-pulse'' sends $\ket0$ to $\tfrac1{\sqrt2}(\ket0-\ii\ket1)$, an equal superposition.
\end{steps}

\begin{example}[Designing a $\pi$-pulse]\label{ex:pi-pulse}
  A qubit has Larmor frequency $\omega_0/2\pi=5$ GHz, and the control field gives a Rabi frequency $\Omega/2\pi=10$ MHz. How long must the field be applied, on resonance, to flip $\ket0$ to $\ket1$? What is the largest flip probability if the field is instead tuned to $5.01$ GHz?
\end{example}
\begin{solution}
  \begin{steps}
    \item \emph{Duration.} We need $\Omega t=\pi$, so $t=\pi/\Omega=\pi/(2\pi\cdot10^7\ \mathrm{s}^{-1})=50$ ns.
    \item \emph{Detuning.} $\Delta/2\pi=5\ \mathrm{GHz}-5.01\ \mathrm{GHz}=-10$ MHz, so $\abs\Delta=\Omega$.
    \item \emph{Largest flip probability.} $\Omega^2/(\Omega^2+\Delta^2)=\tfrac12$: a mistuning of $0.2\%$ in frequency makes a perfect flip impossible.
  \end{steps}
  Numbers of this kind (gigahertz levels, nanosecond pulses) are typical of superconducting qubits.
\end{solution}

\subsection{The quantum Zeno effect}

Measurements change states, and frequent measurements can stop a state from changing at all. Take the Hamiltonian $H=\hbar\omega\sx$ of Example~\ref{ex:H-eq-X}: starting from $\ket0$, the state is $\cos(\omega t)\ket0-\ii\sin(\omega t)\ket1$, and after the time $T=\pi/2\omega$ it is $\ket1$ (up to a phase) with certainty. Now measure in the basis $\ket0,\ket1$ at the times $T/N,2T/N,\dots,T$.

\begin{proposition}[Quantum Zeno effect]\label{prop:zeno}
  The probability that all $N$ measurements give $\ket0$ is $\cos^{2N}(\omega T/N)$, which tends to $1$ as $N\to\infty$.
\end{proposition}
\begin{proof}
  \begin{steps}
    \item \emph{One interval.} Starting from $\ket0$, after time $T/N$ the probability of $\ket0$ is $\cos^2(\omega T/N)$, and if the outcome is $\ket0$ the state is reset to $\ket0$ (Postulate~\ref{post:measurement}).
    \item \emph{All intervals.} Each interval starts again from $\ket0$, so the probabilities multiply: $\cos^{2N}(\omega T/N)$.
    \item \emph{Limit.} With $x=\omega T/N$, $\cos^2x=1-\sin^2x\ge1-x^2$, so $\cos^{2N}(\omega T/N)\ge\bigl(1-\omega^2T^2/N^2\bigr)^N\ge1-\omega^2T^2/N$ by Bernoulli's inequality $(1-a)^N\ge1-Na$. The bound tends to $1$. \qedhere
  \end{steps}
\end{proof}

With $\omega T=\pi/2$, ten measurements already keep the state in $\ket0$ with probability $\cos^{20}(\pi/20)\approx0.78$, and a hundred with probability $0.98$, whereas without measurements the state would certainly have left $\ket0$. The reason is that for short times a state leaves its initial value only quadratically, $1-\cos^2(\omega t)\approx\omega^2t^2$, and each measurement resets the clock. The effect has been observed with trapped ions, and it is one reason why an unintended ``measurement'' by the environment, the decoherence of Section~\ref{sec:interference}, can freeze or spoil a quantum computation.

\begin{exerciseset}
  \exercise{1}\label{ex:spin-eig} Show that for a unit vector $\vec n$, $(\vec n\cdot\vec\sigma)^2=\Id$ (use Proposition~\ref{prop:pauli}), and deduce that the eigenvalues of $\vec n\cdot\vec\sigma$ are $\pm1$. Find the eigenvectors for $\vec n=(1,0,0)$ and $\vec n=(0,1,0)$.

  \exercise{1} In Figure~\ref{fig:sg}, replace the middle SG$_x$ by an apparatus oriented at angle $\theta$ from the $z$-axis in the $xz$-plane. Show that the fraction of atoms that reaches the lower final spot is $\tfrac12\cos^2\tfrac\theta2\sin^2\tfrac\theta2$, and find the angle that maximises it.

  \exercise{2}\label{ex:rotate-x} Let $R(t)=e^{-\ii\omega t\sz/2}=\diag(e^{-\ii\omega t/2},e^{\ii\omega t/2})$. Compute $R(t)\,\sx\,R(t)^\dagger$ as a matrix and show that it equals $\cos(\omega t)\sx+\sin(\omega t)\sy$.

  \exercise{2}\label{ex:precess-n} For $H=\tfrac{\hbar\omega_0}2\,\vec n\cdot\vec\sigma$ with $\abs{\vec n}=1$, show using Lemma~\ref{lem:heisenberg} and $\comm{\sigma_j}{\sigma_k}=2\ii\sum_l\varepsilon_{jkl}\sigma_l$ that the Bloch vector obeys $\dot{\vec r}=\omega_0\,\vec n\times\vec r$. Deduce that $\vec n\cdot\vec r$ and $\abs{\vec r}$ are constant.

  \exercise{2} On resonance, find the duration of a $\pi/2$-pulse for the qubit of Example~\ref{ex:pi-pulse}, and the state, in the rotating frame, after two successive $\pi/2$-pulses.

  \exercise{2} In the Zeno set-up, how many measurements are needed for the probability of staying in $\ket0$ to exceed $0.99$? (Use the bound in the proof, then check numerically.)

  \exercise{3} Show that the Rabi formula is consistent with Lemma~\ref{lem:heisenberg}: in the rotating frame, compute $\tfrac{d}{dt}\langle\sz\rangle$ at $t=0$ and $\tfrac{d^2}{dt^2}\langle\sz\rangle$ at $t=0$, and compare with the Taylor expansion of $1-2p_1(t)$.
\end{exerciseset}

%======================================================================
\section{Thermal states and black-body radiation}\label{sec:blackbody}
%======================================================================

\begin{objectives}
  You will show that the state of largest entropy at a given average energy is the thermal (Gibbs) state $e^{-\beta H}/Z$, compute the average energy of a mode of light whose levels are spaced by $h\nu$, and derive Planck's law, which resolves the ultraviolet catastrophe of the Introduction.
\end{objectives}

\begin{recall}
  The von Neumann entropy (Section~\ref{sec:vonneumann}); the quantum relative entropy and Klein's inequality (Section~\ref{sec:relent}); the functional calculus (Section~\ref{sec:operators}); energy levels (Section~\ref{sec:levels}).
\end{recall}

The black-body spectrum of Figure~\ref{fig:blackbody} is a statement about light in thermal equilibrium with matter at a temperature $T$. To explain it we need two things: the state of a system in thermal equilibrium, which is a question about entropy, and the energy levels of light, which is where Planck's hypothesis $E=h\nu$ enters. Throughout this section entropies are measured with the natural logarithm, $\Ss(\rho)=-\Tr(\rho\ln\rho)$; this only multiplies the entropies of Chapter~\ref{ch:info} by the constant $\ln2$, and changes none of their properties.

\subsection{The state of largest entropy}

A system in contact with a large body at temperature $T$, such as the walls of an oven, exchanges energy with it until it reaches equilibrium. It is a basic principle of statistical physics that the equilibrium state is the one that is \emph{least informative}, that is, of largest entropy, among the states compatible with what is fixed, here the average energy. The relative entropy identifies that state.

\begin{theorem}[Maximum entropy at fixed energy]\label{thm:maxent}
  Let $H$ be a Hamiltonian on a finite-dimensional space, $\beta\in\R$, and
  \[
    \rho_\beta=\frac{e^{-\beta H}}{Z},\qquad Z=\Tr e^{-\beta H}.
  \]
  If $\rho$ is any density operator with the same average energy, $\Tr(\rho H)=\Tr(\rho_\beta H)$, then $\Ss(\rho)\le\Ss(\rho_\beta)$, with equality only if $\rho=\rho_\beta$.
\end{theorem}
\begin{proof}
  \begin{steps}
    \item \emph{$\rho_\beta$ is a state.} If $H=\sum_nE_n\ketbra{E_n}{E_n}$, then $e^{-\beta H}=\sum_ne^{-\beta E_n}\ketbra{E_n}{E_n}$ is positive with trace $Z$, so $\rho_\beta$ is positive with trace $1$. It is invertible, so $\supp\rho\subseteq\supp\rho_\beta$ for every $\rho$.
    \item \emph{Its logarithm.} By the functional calculus, $\ln\rho_\beta=-\beta H-(\ln Z)\Id$.
    \item \emph{Relative entropy with respect to $\rho_\beta$.} For any $\rho$,
      \[
        \Srel\rho{\rho_\beta}=\Tr(\rho\ln\rho)-\Tr(\rho\ln\rho_\beta)=-\Ss(\rho)+\beta\Tr(\rho H)+\ln Z .
      \]
    \item \emph{The same formula for $\rho=\rho_\beta$} gives $0=-\Ss(\rho_\beta)+\beta\Tr(\rho_\beta H)+\ln Z$.
    \item \emph{Subtract,} using $\Tr(\rho H)=\Tr(\rho_\beta H)$: $\Ss(\rho_\beta)-\Ss(\rho)=\Srel\rho{\rho_\beta}$.
    \item \emph{Klein's inequality} (Theorem~\ref{thm:klein}) says that the right side is $\ge0$, with equality only if $\rho=\rho_\beta$. \qedhere
  \end{steps}
\end{proof}

The state $\rho_\beta$ is the \term{thermal state}, or Gibbs state, and $Z$ is the \term{partition function}. Thermodynamics identifies the parameter with the temperature, $\beta=1/k_BT$, where $k_B\approx1.381\times10^{-23}$ J/K is Boltzmann's constant; we take this as the definition of temperature. Three remarks:
\begin{steps}
  \item $\rho_\beta$ is diagonal in the energy basis, and the level $E_n$ has probability $e^{-\beta E_n}/Z$: higher levels are exponentially less likely, the more so the colder the system.
  \item At infinite temperature ($\beta=0$) the thermal state is $\Id/d$, of maximal entropy $\ln d$ (Proposition~\ref{prop:vn-props}); as $T\to0$ it concentrates on the ground level.
  \item The average energy is a logarithmic derivative of $Z$: since $\frac{d}{d\beta}Z=-\sum_nE_ne^{-\beta E_n}$,
    \begin{equation}\label{eq:mean-energy}
      \langle H\rangle_\beta=\Tr(\rho_\beta H)=\frac{\sum_nE_ne^{-\beta E_n}}{\sum_ne^{-\beta E_n}}=-\frac{d}{d\beta}\ln Z .
    \end{equation}
\end{steps}

\begin{example}[A thermal qubit]\label{ex:thermal-qubit}
  A qubit has levels $0$ and $\epsilon>0$. Find its thermal state and its average energy, and their limits as $T\to0$ and $T\to\infty$.
\end{example}
\begin{solution}
  \begin{steps}
    \item \emph{Partition function.} $Z=1+e^{-\beta\epsilon}$.
    \item \emph{Thermal state.} $\rho_\beta=\diag\bigl(1,e^{-\beta\epsilon}\bigr)/(1+e^{-\beta\epsilon})$ in the energy basis.
    \item \emph{Average energy.} By \eqref{eq:mean-energy}, $\langle H\rangle_\beta=\dfrac{\epsilon\,e^{-\beta\epsilon}}{1+e^{-\beta\epsilon}}=\dfrac{\epsilon}{e^{\beta\epsilon}+1}$.
    \item \emph{Cold.} As $T\to0$ ($\beta\to\infty$), $e^{-\beta\epsilon}\to0$, so $\rho_\beta\to\ketbra00$, the ground state, and $\langle H\rangle_\beta\to0$.
    \item \emph{Hot.} As $T\to\infty$ ($\beta\to0$), $e^{-\beta\epsilon}\to1$, so $\rho_\beta\to\tfrac12\Id$ and $\langle H\rangle_\beta\to\epsilon/2$: the two levels become equally likely, and no temperature makes the excited level the more likely one.
  \end{steps}
\end{solution}

\subsection{A mode of light}

Light in a cavity, such as the inside of an oven, is a superposition of standing waves, the \term{modes} of the cavity, each with its own frequency $\nu$. Classically each mode behaves as an oscillator and can hold any amount of energy. Quantum mechanically a mode of frequency $\nu$ is a system with the levels
\[
  E_n=nh\nu,\qquad n=0,1,2,\dots,
\]
and the state of energy $E_n$ is said to contain $n$ \emph{photons} of frequency $\nu$. This is Planck's hypothesis, and it is a theorem of quantum mechanics: the levels of a harmonic oscillator of frequency $\nu$ are $(n+\tfrac12)h\nu$, equally spaced as for no other system, and the constant shift $\tfrac12h\nu$ does not affect anything below.

The space of a mode is infinite-dimensional, but the thermal state is diagonal in the basis of levels, so everything we need is a statement about the probabilities $p_n=e^{-\beta E_n}/Z$, and the sums that appear converge.

\begin{proposition}[Average energy of a mode]\label{prop:planck-mode}
  In the thermal state at temperature $T$, a mode of frequency $\nu$ has average energy and average number of photons
  \[
    \langle E\rangle=\frac{h\nu}{e^{h\nu/k_BT}-1},\qquad\bar n=\frac{\langle E\rangle}{h\nu}=\frac1{e^{h\nu/k_BT}-1}.
  \]
\end{proposition}
\begin{proof}
  Write $x=\beta h\nu$.
  \begin{steps}
    \item \emph{Partition function} (a geometric series, since $e^{-x}<1$): $Z=\sum_{n\ge0}e^{-nx}=\dfrac1{1-e^{-x}}$.
    \item \emph{Logarithm.} $\ln Z=-\ln\bigl(1-e^{-\beta h\nu}\bigr)$.
    \item \emph{Differentiate} as in \eqref{eq:mean-energy}: $\langle E\rangle=-\dfrac{d}{d\beta}\ln Z=\dfrac{h\nu\,e^{-\beta h\nu}}{1-e^{-\beta h\nu}}=\dfrac{h\nu}{e^{\beta h\nu}-1}$. \qedhere
  \end{steps}
\end{proof}

Compare with the classical answer, $\langle E\rangle=k_BT$ for every mode (the ``equipartition of energy'').
\begin{steps}
  \item \emph{Low frequency,} $h\nu\ll k_BT$: $e^x-1\approx x$, so $\langle E\rangle\approx h\nu/x=k_BT$. Many photons are present, the discreteness of the levels is invisible, and the classical answer is recovered.
  \item \emph{High frequency,} $h\nu\gg k_BT$: $\langle E\rangle\approx h\nu\,e^{-h\nu/k_BT}\to0$. The first excited level costs far more than the typical thermal energy $k_BT$, so it is almost never occupied: the mode is \emph{frozen}.
\end{steps}
It is this freezing of the high-frequency modes that saves the total energy from being infinite.

\subsection{Planck's law}

The number of modes of a cavity of volume $V$ with frequencies between $\nu$ and $\nu+d\nu$ is $8\pi V\nu^2d\nu/c^3$, where $c$ is the speed of light (this is a count of standing waves, Exercise~\ref{ex:mode-count}). Multiplying by the energy per mode, the energy per unit volume and unit frequency is
\[
  u(\nu)=\frac{8\pi\nu^2}{c^3}\,\langle E\rangle=\frac{8\pi h\nu^3}{c^3}\,\frac1{e^{h\nu/k_BT}-1}\qquad\text{(Planck)},
\]
against $u_{\mathrm{cl}}(\nu)=8\pi\nu^2k_BT/c^3$ with the classical $\langle E\rangle=k_BT$ (Rayleigh--Jeans). In the variable $x=h\nu/k_BT$ the two are proportional to $x^3/(e^x-1)$ and $x^2$, the two curves of Figure~\ref{fig:blackbody}. The classical one grows without bound and its integral over all frequencies diverges: the ultraviolet catastrophe. Planck's law agrees with it at low frequencies, where $x^3/(e^x-1)=x^2(1+O(x))$, and falls exponentially at high ones. Planck found this formula in 1900~\cite{planck1901} by fitting measurements; we have derived it from the discreteness of the levels of a mode and the principle of maximal entropy.

\begin{example}[The peak of the spectrum and the cosmic background]\label{ex:wien}
  Show that $x^3/(e^x-1)$ has its maximum where $3(1-e^{-x})=x$, at $x_*\approx2.82$. Deduce that the peak frequency of a black body is proportional to its temperature, and compute it for the cosmic microwave background, the radiation left over from the early universe, whose temperature is $T\approx2.725$ K. Use $h\approx6.626\times10^{-34}$ J\,s.
\end{example}
\begin{solution}
  \begin{steps}
    \item \emph{Differentiate.} $\dfrac{d}{dx}\dfrac{x^3}{e^x-1}=\dfrac{3x^2(e^x-1)-x^3e^x}{(e^x-1)^2}$, which vanishes for $x>0$ when $3(e^x-1)=xe^x$, that is, $3(1-e^{-x})=x$.
    \item \emph{Solve numerically.} The function $g(x)=3(1-e^{-x})-x$ has $g(2.8)\approx0.018>0$ and $g(2.85)\approx-0.024<0$; bisection gives $x_*\approx2.821$.
    \item \emph{Wien's law.} The peak is at $h\nu_{\max}=x_*k_BT$, so $\nu_{\max}=2.82\,k_BT/h$ is proportional to $T$: hotter bodies glow at higher frequencies, from red to white.
    \item \emph{The cosmic background.} $\nu_{\max}\approx2.82\cdot1.381\times10^{-23}\cdot2.725/6.626\times10^{-34}\approx1.60\times10^{11}$ Hz, that is, $160$ GHz, in the microwave range.
  \end{steps}
  The cosmic microwave background is the most perfect black body ever measured: its spectrum follows Planck's law to better than one part in ten thousand.
\end{solution}

\begin{lookahead}
  Thermal states are the natural noise of the physical world. A qubit at temperature $T$ is in the state of Example~\ref{ex:thermal-qubit}, partly excited, unless its level spacing $\epsilon$ is much larger than $k_BT$. For superconducting qubits, with $\epsilon/h\approx5$ GHz, this requires temperatures of about $10$ mK, which is why they are operated in dilution refrigerators.
\end{lookahead}

\begin{exerciseset}
  \exercise{1} Compute the thermal state and the average energy of a qubit with $\epsilon/h=5$ GHz at $T=20$ mK and at $T=300$ K. (Use $h/k_B\approx4.8\times10^{-11}$ K\,s.)

  \exercise{1} Check that $\rho_\beta$ commutes with $H$, and deduce that the thermal state is stationary: $e^{-\ii Ht/\hbar}\rho_\beta e^{\ii Ht/\hbar}=\rho_\beta$.

  \exercise{2} Show that the entropy of the thermal state of a mode, in terms of the average number of photons $\bar n$, is $\Ss=(\bar n+1)\ln(\bar n+1)-\bar n\ln\bar n$. (Hint: $p_n=\bar n^n/(\bar n+1)^{n+1}$.)

  \exercise{2} Show that $\Var(E)=\dfrac{d^2}{d\beta^2}\ln Z$ for a thermal state, and compute it for a mode. Deduce $\Var(n)=\bar n^2+\bar n$, which is larger than for a Poisson distribution with the same mean.

  \exercise{2} Using $\int_0^\infty\frac{x^3}{e^x-1}\dd x=\frac{\pi^4}{15}$, show that the total energy per unit volume of black-body radiation is proportional to $T^4$ (the Stefan--Boltzmann law). Why is the classical answer infinite?

  \exercise{3}\label{ex:mode-count} The standing waves in a cubic cavity of side $L$ have wave vectors $\vec k=\tfrac\pi L(j_1,j_2,j_3)$ with positive integers $j_i$, frequency $\nu=c\abs{\vec k}/2\pi$, and two polarisations each. Show that the number of modes with frequency below $\nu$ is approximately $\frac{8\pi}{3}\frac{L^3\nu^3}{c^3}$ for large $\nu$, by counting points of the integer lattice in an eighth of a ball, and deduce the mode density $8\pi V\nu^2/c^3$.

  \exercise{3} Theorem~\ref{thm:maxent} fixes the average energy. Show in the same way that, if nothing is fixed, the state of largest entropy is $\Id/d$, and that, if the average of each of several commuting observables $A_1,\dots,A_k$ is fixed, the maximiser has the form $e^{-\sum_j\beta_jA_j}/Z$ whenever a state of that form with the right averages exists.
\end{exerciseset}

%======================================================================
\section{No cloning, and teleportation}\label{sec:teleportation}
%======================================================================

\begin{objectives}
  You will prove that no physical process can copy an unknown quantum state, and see how an unknown state can nevertheless be transferred from one place to another, destroying the original, with one shared entangled pair and two classical bits.
\end{objectives}

\begin{recall}
  Unitaries preserve inner products (Section~\ref{sec:operators}); non-orthogonal states cannot be distinguished (Theorem~\ref{thm:no-distinguish}); the Bell states, the controlled-NOT and superdense coding (Sections~\ref{sec:composite} and~\ref{sec:superdense}); the reduced state of half a Bell pair (Example~\ref{ex:reduced-bell}).
\end{recall}

The last two phenomena are not old puzzles but consequences of the postulates that were noticed only in the 1980s and 1990s. They are at the heart of quantum information: the first one makes quantum cryptography possible, and the second one shows that entanglement is a resource that can be spent.

\subsection{Unknown states cannot be copied}

Classical information can be copied at will. A copying machine for qubits would take a qubit in an unknown state $\ket\psi$ and a second qubit in a fixed ``blank'' state $\ket e$, and produce two copies of $\ket\psi$. By the postulates it would act by a unitary on the pair. Wootters and Zurek~\cite{wootters1982}, and independently Dieks~\cite{dieks1982}, showed that no such machine exists.

\begin{theorem}[No-cloning, Wootters--Zurek and Dieks, 1982]\label{thm:no-cloning}
  Let $U$ be a unitary on $\Hil\tp\Hil$ and $\ket e\in\Hil$ a unit vector. If $U\bigl(\ket\psi\ket e\bigr)=\ket\psi\ket\psi$ and $U\bigl(\ket\varphi\ket e\bigr)=\ket\varphi\ket\varphi$ for two states $\ket\psi,\ket\varphi$, then either $\braket\varphi\psi=0$ or $\ket\varphi=\ket\psi$ up to a phase. In particular, no unitary copies every state.
\end{theorem}
\begin{proof}
  \begin{steps}
    \item \emph{Inner product before.} $\bigl(\bra\varphi\bra e\bigr)\bigl(\ket\psi\ket e\bigr)=\braket\varphi\psi\braket ee=\braket\varphi\psi$ (Proposition~\ref{prop:tensor-inner}).
    \item \emph{Inner product after.} $\bigl(\bra\varphi\bra\varphi\bigr)\bigl(\ket\psi\ket\psi\bigr)=\braket\varphi\psi^2$.
    \item \emph{Unitaries preserve inner products,} so $s=\braket\varphi\psi$ satisfies $s=s^2$.
    \item \emph{Solve.} $s=0$ or $s=1$. In the second case $\abs{\braket\varphi\psi}=1$, which is the equality case of Cauchy--Schwarz (Theorem~\ref{thm:cs}): $\ket\varphi$ is a multiple of $\ket\psi$, and the multiple has modulus $1$. \qedhere
  \end{steps}
\end{proof}

A machine can therefore copy the states of one orthonormal basis, and no others. The controlled-NOT does exactly that.

\begin{example}[The CNOT copies bits but not qubits]\label{ex:cnot-clone}
  Compute $\CNOT\bigl(\ket x\ket0\bigr)$ for $x=0,1$, and $\CNOT\bigl(\ket+\ket0\bigr)$. Is the latter $\ket+\ket+$?
\end{example}
\begin{solution}
  \begin{steps}
    \item \emph{Basis states.} $\CNOT\ket0\ket0=\ket0\ket0$ and $\CNOT\ket1\ket0=\ket1\ket1$: both are copied.
    \item \emph{The superposition.} By linearity, $\CNOT\ket+\ket0=\tfrac1{\sqrt2}\bigl(\ket{00}+\ket{11}\bigr)=\ket{\Phi^+}$.
    \item \emph{Compare.} $\ket+\ket+=\tfrac12\bigl(\ket{00}+\ket{01}+\ket{10}+\ket{11}\bigr)\neq\ket{\Phi^+}$; indeed $\ket{\Phi^+}$ is entangled, so it is not of the form $\ket a\ket b$ at all.
  \end{steps}
  Linearity forces the copy of a superposition to be the superposition of the copies, and that is an entangled state, not two copies.
\end{solution}

The argument does not depend on the copier being a unitary on two qubits only: it may use an auxiliary system of any size, which ends in a state that depends on $\ket\psi$ (Exercise~\ref{ex:clone-ancilla}). Nor can a measurement do the job, since the statistics of one measurement on one copy do not determine $\ket\psi$. No-cloning is consistent with everything else we know.
\begin{steps}
  \item If cloning were possible, we could make many copies of a state that is either $\ket0$ or $\ket+$, measure them, and decide which state it was with as small an error as we like, contradicting Theorem~\ref{thm:no-distinguish}.
  \item If cloning were possible, Bob could copy his half of a singlet many times and learn which measurement Alice had made on her half, and Alice could signal to him faster than light, which Exercise~\ref{ex:no-signalling} forbids.
\end{steps}

\begin{lookahead}
  No-cloning is the reason why quantum key distribution is secure: an eavesdropper cannot copy the qubits sent between Alice and Bob in order to measure the copies later, and any measurement she makes on the originals disturbs them in a way that Alice and Bob can detect. It is also why quantum error correction cannot simply keep backup copies of qubits, and must instead spread their information over entangled states of several qubits.
\end{lookahead}

\subsection{Teleportation}

Alice holds a qubit $C$ in a state $\ket\psi=a\ket0+b\ket1$ that she does not know, and she wants Bob, far away, to have it. She cannot measure $\ket\psi$ and send a description of it, since one measurement does not determine $a$ and $b$; she cannot copy it; and perhaps she cannot send the qubit itself. If, however, Alice and Bob share in advance a Bell pair $\bell00$, of which Alice holds qubit $A$ and Bob qubit $B$, then two classical bits are enough. The protocol is due to Bennett, Brassard, Cr\'epeau, Jozsa, Peres and Wootters~\cite{bennett1993}.

\paragraph{The protocol.}
\begin{stepsn}
  \item Alice measures her two qubits $C$ and $A$ in the Bell basis $\{\bell xy\}$, and obtains two bits $x,y$.
  \item She sends $x,y$ to Bob over an ordinary classical channel.
  \item Bob applies $Z^xX^y$ to his qubit $B$. It is now in the state $\ket\psi$.
\end{stepsn}

\begin{figure}[htb]
  \centering
  \begin{tikzpicture}[>=stealth,
      box/.style={draw=navy,thick,fill=navy!8,minimum height=0.8cm,font=\small},
      cl/.style={draw=figacc,thick,fill=figacc!10,minimum height=0.8cm,font=\small}]
    \node[font=\small,anchor=east] at (-0.1,1.2) {$C$: $\ket\psi$};
    \node[draw=ink,thick,fill=ink!8,minimum width=1.6cm,minimum height=1.6cm,font=\small,align=center] (src) at (1.2,-0.6) {$\bell00$};
    \node[font=\footnotesize,ink,anchor=north] at (1.2,-1.45) {shared pair};
    \node[cl,minimum width=2.2cm,minimum height=2.0cm,align=center] (bm) at (5.0,0.6) {Bell\\measurement};
    \node[font=\footnotesize,ink] at (5.0,1.95) {Alice};
    \node[box,minimum width=1.4cm] (bob) at (9.2,-1.0) {$Z^xX^y$};
    \node[font=\footnotesize,ink] at (9.2,-1.75) {Bob};
    \draw[thick] (0,1.2) -- (bm.west|-0,1.2);
    \draw[thick] (src.east|-0,0) -- node[above,font=\footnotesize] {$A$} (bm.west|-0,0);
    \draw[thick] (src.east|-0,-1.0) -- node[above,font=\footnotesize] {$B$} (bob.west);
    \draw[->,thick] (bob.east) -- ++(1.0,0) node[right,font=\small] {$\ket\psi$};
    \draw[figacc,double,thick,->] (bm.east) -- ++(1.6,0) -| node[pos=0.75,right,font=\footnotesize] {bits $x,y$} (bob.north);
  \end{tikzpicture}
  \caption{Teleportation. Alice performs a Bell measurement on the unknown qubit $C$ and her half $A$ of a shared Bell pair, and sends the two resulting bits (double line) to Bob, whose correction $Z^xX^y$ turns his half $B$ into $\ket\psi$.}
  \label{fig:teleport}
\end{figure}

Everything rests on one identity, which rewrites the initial state of the three qubits in the Bell basis of Alice's two.

\begin{lemma}\label{lem:teleport-identity}
  With the qubits in the order $C,A,B$,
  \[
    \ket\psi_C\,\bell00_{AB}=\frac12\sum_{x,y\in\{0,1\}}\bell xy_{CA}\;X^yZ^x\ket\psi_B .
  \]
\end{lemma}
\begin{proof}
  \begin{steps}
    \item \emph{Expand.} $\ket\psi\bell00=\tfrac1{\sqrt2}\bigl(a\ket{00}\ket0+a\ket{01}\ket1+b\ket{10}\ket0+b\ket{11}\ket1\bigr)$, where the two-qubit kets refer to $C,A$ and the last one to $B$.
    \item \emph{Invert the Bell states.} From the definitions of Section~\ref{sec:composite}, $\ket{00}=\tfrac1{\sqrt2}(\bell00+\bell10)$, $\ket{11}=\tfrac1{\sqrt2}(\bell00-\bell10)$, $\ket{01}=\tfrac1{\sqrt2}(\bell01+\bell11)$ and $\ket{10}=\tfrac1{\sqrt2}(\bell01-\bell11)$.
    \item \emph{Substitute and collect} the terms of each Bell state:
      \[
        \begin{aligned}
          \ket\psi\bell00=\tfrac12\Bigl(&\bell00\bigl(a\ket0+b\ket1\bigr)+\bell10\bigl(a\ket0-b\ket1\bigr)\\
          &+\bell01\bigl(b\ket0+a\ket1\bigr)+\bell11\bigl(-b\ket0+a\ket1\bigr)\Bigr).
        \end{aligned}
      \]
    \item \emph{Recognise Bob's states.} $a\ket0-b\ket1=Z\ket\psi$, $b\ket0+a\ket1=X\ket\psi$, and $-b\ket0+a\ket1=X(a\ket0-b\ket1)=XZ\ket\psi$. \qedhere
  \end{steps}
\end{proof}

\begin{theorem}[Teleportation]\label{thm:teleportation}
  In the protocol above, each pair of bits $x,y$ occurs with probability $\tfrac14$, and after Bob's correction his qubit is in the state $\ket\psi$, whatever $\ket\psi$ was.
\end{theorem}
\begin{proof}
  \begin{steps}
    \item \emph{Alice's measurement} has the projectors $\ketbra{\beta_{xy}}{\beta_{xy}}\tp\Id_B$, since the Bell states are an orthonormal basis of $C\tp A$.
    \item \emph{Probabilities.} By Lemma~\ref{lem:teleport-identity}, the component of the state along $\bell xy$ is $\tfrac12\bell xy\,X^yZ^x\ket\psi$, of squared norm $\tfrac14$, since $X^yZ^x$ is unitary.
    \item \emph{Post-measurement state.} After outcome $x,y$ the three qubits are in $\bell xy_{CA}\,X^yZ^x\ket\psi_B$ (Postulate~\ref{post:measurement}): Bob's qubit is in $X^yZ^x\ket\psi$.
    \item \emph{Correction.} $Z^xX^y\cdot X^yZ^x=Z^x\,(X^y)^2\,Z^x=\Id$, because $X^2=Z^2=\Id$. So Bob ends with $\ket\psi$. \qedhere
  \end{steps}
\end{proof}

The protocol obeys every rule we have met.
\begin{steps}
  \item \emph{No cloning.} Alice's qubit $C$ ends in a Bell state with $A$, so no trace of $\ket\psi$ is left with her: the state has been moved, not copied.
  \item \emph{No faster-than-light signalling.} Until the two bits arrive, Bob does not know which correction to apply, and his qubit carries no information about $\ket\psi$ (Example~\ref{ex:bob-before}). The bits travel no faster than light.
  \item \emph{An exchange rate.} One shared Bell pair and two classical bits transmit one qubit. Superdense coding (Section~\ref{sec:superdense}) is the reverse exchange: one shared Bell pair and one qubit transmit two classical bits. Neither protocol works without the shared entanglement, which is consumed in the process.
  \item \emph{Nothing is learnt.} The bits $x,y$ are uniformly random whatever $\ket\psi$ is, so they reveal nothing about $a$ and $b$. All the information about $\ket\psi$ reaches Bob through the entanglement.
\end{steps}

\begin{example}[Bob's qubit before the bits arrive]\label{ex:bob-before}
  Show that, after Alice's measurement but before Bob learns its outcome, Bob's qubit is in the state $\tfrac12\Id$, whatever $\ket\psi$ is.
\end{example}
\begin{solution}
  \begin{steps}
    \item \emph{Bob's description.} Not knowing $x,y$, Bob describes his qubit by the ensemble of the four states $X^yZ^x\ket\psi$ with probabilities $\tfrac14$ (Section~\ref{sec:density}):
      \[
        \rho_B=\tfrac14\sum_{x,y}X^yZ^x\ketbra\psi\psi Z^xX^y .
      \]
    \item \emph{Use the Bloch form.} $\ketbra\psi\psi=\tfrac12(\Id+r_x\sx+r_y\sy+r_z\sz)$ (Proposition~\ref{prop:bloch}).
    \item \emph{Average over $x$.} Conjugating by $Z$ keeps $\Id$ and $\sz$ and changes the signs of $\sx,\sy$ (they anticommute with $\sz$). The average over $x=0,1$ kills the $\sx$ and $\sy$ terms, and leaves $\tfrac12(\Id+r_z\sz)$.
    \item \emph{Average over $y$.} Conjugating by $X$ changes the sign of $\sz$, and the average over $y$ leaves $\tfrac12\Id$.
  \end{steps}
  This is also clear from Section~\ref{sec:density}: Alice's operations act on her qubits only, so they cannot change Bob's reduced state, which was $\tfrac12\Id$ from the start (Example~\ref{ex:reduced-bell}).
\end{solution}

\begin{lookahead}
  Teleportation was first realised with photons in 1997, and has since been performed between islands and between the ground and a satellite more than a thousand kilometres away. In quantum computing it is a working tool: qubits are teleported between distant parts of a processor or of a network, and gates can be teleported too, which is a key idea in fault-tolerant quantum computation.
\end{lookahead}

\begin{exerciseset}
  \exercise{1} Verify Lemma~\ref{lem:teleport-identity} for $\ket\psi=\ket1$ by expanding both sides in the computational basis.

  \exercise{1} Run the protocol for $\ket\psi=\ket+$ and the outcome $x=1$, $y=0$: write the state of Bob's qubit before and after his correction.

  \exercise{2}\label{ex:clone-ancilla} Suppose a cloner uses an auxiliary system, $U\bigl(\ket\psi\ket e\ket m\bigr)=\ket\psi\ket\psi\ket{m_\psi}$, where the final auxiliary state $\ket{m_\psi}$ may depend on $\ket\psi$. Show that $\abs{\braket\varphi\psi}\le\abs{\braket\varphi\psi}^2$ for any two states it clones, and conclude as in Theorem~\ref{thm:no-cloning}.

  \exercise{2} Alice and Bob share $\bell11$ instead of $\bell00$. Rewrite Lemma~\ref{lem:teleport-identity} for this case, and find the correction Bob must apply for each outcome $x,y$.

  \exercise{2} Show that a cloner of the two states $\ket0$ and $\ket+$ would allow us to distinguish them with error probability as small as we like, and explain why this contradicts Theorem~\ref{thm:no-distinguish}.

  \exercise{3} \emph{Entanglement swapping.} Alice and Bob share a Bell pair $\bell00_{AB}$, and Bob and Charlie share another, $\bell00_{B'C}$. Bob teleports his qubit $B$ to Charlie using the second pair. Show that afterwards Alice's qubit and Charlie's qubit are in a Bell state (which one, after Charlie's correction?), although they have never interacted.

  \exercise{3} An imperfect copier sends $\ket0\ket e\mapsto\ket0\ket0$ and $\ket1\ket e\mapsto\ket1\ket1$ (the CNOT of Example~\ref{ex:cnot-clone}). Show that for $\ket\psi=\cos\tfrac\theta2\ket0+e^{\ii\phi}\sin\tfrac\theta2\ket1$ the reduced state of each output qubit is $\cos^2\tfrac\theta2\ketbra00+\sin^2\tfrac\theta2\ketbra11$, and that its fidelity $\bra\psi\rho\ket\psi$ with $\ket\psi$ is $\tfrac12(1+\cos^2\theta)$. Which input states are copied worst?
\end{exerciseset}

\begin{furtherreading}
  \item Feynman, Leighton and Sands~\cite{feynman1965}: the classic physical introduction to quantum mechanics, which starts, as this chapter does, from interference with two paths and from two-level systems.
  \item Englert~\cite{englert1996}: the sharp form of the trade-off between fringe visibility and which-path information of Section~\ref{sec:interference}; Schlosshauer~\cite{schlosshauer2007}: decoherence in depth.
  \item Griffiths and Schroeter~\cite{griffiths2018} and Sakurai and Napolitano~\cite{sakurai2020}: the exact solution of the hydrogen atom and the harmonic oscillator, spin and magnetic resonance, and the statistical mechanics behind Section~\ref{sec:blackbody}.
  \item Wootters and Zurek~\cite{wootters1982} and Dieks~\cite{dieks1982}: the original no-cloning theorems, each one page long; Bennett et~al.~\cite{bennett1993}: the original teleportation paper. Nielsen and Chuang~\cite{nielsenchuang2010}, Sections~1.3 and~12.1, treat both in the language of circuits.
\end{furtherreading}

